% =====================================================================
% sdt_foundations.tex — companion derivation note for the VDA rebuild.
%
% Writes up the signal-detection foundations of the rebuilt model's
% "inherited skeleton" (Rebuild/manuscript/sections/model.tex, §2.1):
% the per-location equal-variance Gaussian decision, the derivation of
% the hit-rate and false-alarm-rate expressions, the modelling
% conventions that make those expressions well-posed, and the bridge
% from a single location to the multi-location no-false-alarm term
% P_no-fa where assumption A1 (independence) enters.
%
% Self-contained. Build: pdflatex sdt_foundations.tex (twice).
% =====================================================================

\documentclass[11pt]{article}

\usepackage[utf8]{inputenc}
\usepackage[T1]{fontenc}
\usepackage{lmodern}
\usepackage[a4paper, margin=1in]{geometry}
\usepackage{amsmath, amssymb, amsthm}
\usepackage{mathtools}
\usepackage{xcolor}
\usepackage[colorlinks=true, allcolors=blue!45!black]{hyperref}

\newtheorem{lemma}{Lemma}
\newtheorem{proposition}{Proposition}
\theoremstyle{remark}
\newtheorem{remark}{Remark}

% shorthand aligned with sections/model.tex
\newcommand{\HR}{\mathrm{HR}}
\newcommand{\FAR}{\mathrm{FAR}}
\newcommand{\CR}{\mathrm{CR}}
\newcommand{\dprime}{d'}
\newcommand{\PnoFA}{P_{\mathrm{no\text{-}fa}}}
\newcommand{\Nloc}{N}

\title{Signal-detection foundations of the value-directed-attention model\\[2pt]
       \large A derivation note for the rebuilt model section}
\author{Jonathan Morgan \and James Herman}
\date{\today}

\begin{document}
\maketitle

\begin{abstract}
\noindent
This note derives, from first principles, the per-location hit-rate
and false-alarm-rate expressions that the rebuilt model inherits from
the Herman Lab manuscript, and states the modelling conventions that
make them well-posed. The goal is to make every symbol in
Equation~(\ref{eq:sdt}) traceable to an explicit probabilistic
statement about an observer's decision, and to localise the single
place where the per-location independence assumption (A1) enters the
$\Nloc$-location reward, which is the seam the rebuild later cuts along
to introduce the decorrelation channel.
\end{abstract}

% ---------------------------------------------------------------------
\section{The task and the decision}
\label{sec:task}

An observer monitors $\Nloc \ge 2$ spatial locations (in the cued
change-detection paradigm, $\Nloc$ oriented patches). On each trial at
most one location changes. The observer cannot read out the world
directly; at each location it has access only to a scalar
\emph{decision variable} $X_i$ --- a noisy internal summary of the
evidence that location $i$ changed. The observer issues a separate
binary judgement at each location, ``change'' or ``no change.'' The
quantities of interest are the two conditional response rates at each
location: the \emph{hit rate} $\HR_i$ (the probability of reporting
``change'' when location $i$ truly changed) and the \emph{false-alarm
rate} $\FAR_i$ (the probability of reporting ``change'' when it did
not).

We work throughout with the \emph{normative} observer: rather than
modelling any particular subject's idiosyncratic policy, we ask what
an observer that has already tuned its sensitivities and criteria to
maximise expected reward would do. The signal-detection layer below is
the measurement model on which that optimisation is later performed.

% ---------------------------------------------------------------------
\section{The per-location signal-detection model}
\label{sec:model}

\paragraph{Distributions.}
Adopt the equal-variance Gaussian model with unit variance
($\sigma = 1$). At location $i$, the decision variable is drawn from
one of two distributions according to whether a change occurred:
\begin{equation}
  \text{no change } (H_0):\ X_i \sim \mathcal{N}\!\left(-\tfrac{\dprime_i}{2},\, 1\right),
  \qquad
  \text{change } (H_1):\ X_i \sim \mathcal{N}\!\left(+\tfrac{\dprime_i}{2},\, 1\right).
  \label{eq:dists}
\end{equation}
The two means are placed symmetrically about $0$, so that their
separation is exactly $\dprime_i$. This is the definition of
sensitivity:
\begin{equation}
  \dprime_i
  \;:=\;
  \frac{\mathbb{E}[X_i \mid H_1] - \mathbb{E}[X_i \mid H_0]}{\sigma}
  \;=\; \tfrac{\dprime_i}{2} - \left(-\tfrac{\dprime_i}{2}\right),
  \qquad \sigma = 1 .
  \label{eq:dprime-def}
\end{equation}
Larger $\dprime_i$ pushes the two distributions apart and makes change
more discriminable from no-change.

\paragraph{Decision rule.}
The observer reports ``change'' at location $i$ when the evidence
exceeds a threshold, the \emph{criterion} $c_i \in \mathbb{R}$:
\begin{equation}
  \text{respond ``change'' at } i
  \iff
  X_i > c_i .
  \label{eq:rule}
\end{equation}
Because $c_i$ is measured on the same centred axis as
(\ref{eq:dists}), $c_i = 0$ is the unbiased criterion (it cuts the two
distributions symmetrically); $c_i > 0$ is conservative (fewer
``change'' responses), $c_i < 0$ liberal.

\paragraph{Target.}
Both $\HR_i$ and $\FAR_i$ are the \emph{same} event under
(\ref{eq:rule}) --- ``$X_i$ exceeds $c_i$'' --- evaluated under the two
hypotheses of (\ref{eq:dists}). They are therefore tail areas of
Gaussians, computed with the standard-normal cumulative distribution
function $\Phi(z) = \Pr(Z \le z)$, $Z \sim \mathcal{N}(0,1)$.

% ---------------------------------------------------------------------
\section{Deriving the hit and false-alarm rates}
\label{sec:derivation}

We use two elementary facts, stated as lemmas so the derivation has no
hidden steps.

\begin{lemma}[Standardisation]
\label{lem:std}
If $X \sim \mathcal{N}(\mu, \sigma^2)$ then $Z := (X-\mu)/\sigma \sim
\mathcal{N}(0,1)$, and for any $t \in \mathbb{R}$,
$\Pr(X > t) = \Pr\!\big(Z > (t-\mu)/\sigma\big)$.
\end{lemma}

\begin{lemma}[Tail symmetry of the standard normal]
\label{lem:sym}
For all $z \in \mathbb{R}$, $\;1 - \Phi(z) = \Phi(-z)$, because the
standard-normal density is even about $0$.
\end{lemma}

\begin{proposition}[Per-location SDT rates]
\label{prop:rates}
Under the model (\ref{eq:dists})--(\ref{eq:rule}) with $\sigma = 1$,
\begin{equation}
  \HR_i = \Phi\!\left(\tfrac{\dprime_i}{2} - c_i\right),
  \qquad
  \FAR_i = \Phi\!\left(-\tfrac{\dprime_i}{2} - c_i\right).
  \label{eq:sdt}
\end{equation}
\end{proposition}

\begin{proof}
\emph{Hit rate.} Condition on $H_1$, so $X_i \sim
\mathcal{N}(\dprime_i/2, 1)$. Apply Lemma~\ref{lem:std} with
$\mu = \dprime_i/2$, $\sigma = 1$, $t = c_i$, and then
Lemma~\ref{lem:sym}:
\begin{align}
  \HR_i
  &= \Pr(X_i > c_i \mid H_1)
   = \Pr\!\big(Z > c_i - \tfrac{\dprime_i}{2}\big)
   = 1 - \Phi\!\big(c_i - \tfrac{\dprime_i}{2}\big)
   = \Phi\!\big(\tfrac{\dprime_i}{2} - c_i\big).
\end{align}
\emph{False-alarm rate.} Condition on $H_0$, so $X_i \sim
\mathcal{N}(-\dprime_i/2, 1)$. Identically,
\begin{align}
  \FAR_i
  &= \Pr(X_i > c_i \mid H_0)
   = \Pr\!\big(Z > c_i + \tfrac{\dprime_i}{2}\big)
   = 1 - \Phi\!\big(c_i + \tfrac{\dprime_i}{2}\big)
   = \Phi\!\big(-\tfrac{\dprime_i}{2} - c_i\big).
\end{align}
These are (\ref{eq:sdt}).
\end{proof}

% ---------------------------------------------------------------------
\section{Conventions and what they buy}
\label{sec:conventions}

\begin{remark}[The centring is a choice, not a law]
Placing the means at $\mp\dprime_i/2$ fixes the origin of the decision
axis at the midpoint of the two distributions. Any other origin gives
an algebraically equivalent model. For instance, putting noise at $0$
and signal at $\dprime_i$ and measuring a criterion $c_i'$ from the
noise mean yields $\HR_i = \Phi(\dprime_i - c_i')$,
$\FAR_i = \Phi(-c_i')$; the substitution $c_i = c_i' - \dprime_i/2$
recovers (\ref{eq:sdt}). The centred convention is preferred only
because it makes $c_i = 0$ the bias-free criterion, so the sign of
$c_i$ reads directly as conservative/liberal.
\end{remark}

\begin{remark}[Why a threshold on $X_i$ is general]
The rule (\ref{eq:rule}) is a simple threshold on the raw decision
variable, not on a likelihood ratio. Under the equal-variance Gaussian
model the likelihood ratio
$\Lambda(x) = \varphi(x - \dprime_i/2)/\varphi(x + \dprime_i/2)
= \exp(\dprime_i x)$ is strictly increasing in $x$, so
``$\Lambda(X_i) > \tau$'' is equivalent to ``$X_i > c_i$'' for a
suitable $c_i$. By the Neyman--Pearson lemma the threshold rule is
therefore the optimal test, and (\ref{eq:rule}) loses no generality.
This equivalence relies on equal variances: with $\sigma_0 \ne
\sigma_1$ the likelihood ratio is non-monotone, the optimal region is
an interval rather than a half-line, and (\ref{eq:sdt}) no longer
holds. Equal variance is thus a load-bearing modelling assumption, not
a cosmetic one.
\end{remark}

% ---------------------------------------------------------------------
\section{From one location to \texorpdfstring{$\Nloc$}{N}: the no-false-alarm term}
\label{sec:pnofa}

The expected reward couples locations only through the no-change
trials, where the observer is penalised for false-alarming
\emph{anywhere}. Define the per-location correct-rejection probability
\begin{equation}
  \CR_i := 1 - \FAR_i
  = 1 - \Phi\!\left(-\tfrac{\dprime_i}{2} - c_i\right)
  = \Phi\!\left(\tfrac{\dprime_i}{2} + c_i\right)
  =: \Phi(b_i),
  \qquad
  b_i := c_i + \tfrac{\dprime_i}{2},
  \label{eq:cr}
\end{equation}
using Lemma~\ref{lem:sym} again; $b_i$ is the $z$-score the no-change
decision variable must stay below to avoid a false alarm. The
probability that \emph{no} location false-alarms on a no-change trial
is then
\begin{equation}
  \PnoFA
  \;=\;
  \Pr\!\left(\bigcap_{i} \{X_i \le c_i\} \,\Big|\, H_0\ \text{at every location}\right).
  \label{eq:pnofa-joint}
\end{equation}
This is the \emph{only} term in the expected reward that involves more
than one location at once: on a change trial the change is at a single
location, so the change-trial reward is linear in the marginal hit
rates and contains no cross-location product.

\paragraph{Where assumption A1 enters.}
If the per-location decision variables are independent (assumption
A1), the joint event (\ref{eq:pnofa-joint}) factorises into a product
of marginals, and with homogeneous uncued allocation
($\dprime_u, c_u$ shared across the $\Nloc-1$ uncued locations) it
collapses to
\begin{equation}
  \PnoFA^{\text{indep}}
  \;=\;
  \prod_{i}\big(1 - \FAR_i\big)
  \;=\;
  \Phi(b_c)\,\Phi(b_u)^{\Nloc - 1}.
  \label{eq:pnofa-indep}
\end{equation}
Equation~(\ref{eq:pnofa-indep}) is the sole appearance of A1 in the
model. The rebuilt manuscript replaces it with the exact
equicorrelated-Gaussian orthant probability: writing the no-change
variables in one-factor form
$X_i = \mu_i + \sqrt{\rho}\,Z + \sqrt{1-\rho}\,\varepsilon_i$ with
$Z, \varepsilon_i \stackrel{\text{iid}}{\sim} \mathcal{N}(0,1)$ and
$\mu_i = -\dprime_i/2$, the variables are conditionally independent
given $Z$, so multiplying the conditional correct-rejection
probabilities and integrating out $Z$ gives
\begin{equation}
  \PnoFA(\rho)
  \;=\;
  \int_{-\infty}^{\infty}
  \Phi\!\left(\frac{b_c - \sqrt{\rho}\,z}{\sqrt{1-\rho}}\right)
  \Phi\!\left(\frac{b_u - \sqrt{\rho}\,z}{\sqrt{1-\rho}}\right)^{\!\Nloc - 1}
  \varphi(z)\,dz,
  \qquad \rho \in [0,1).
  \label{eq:pnofa-rho}
\end{equation}
At $\rho = 0$ the integrand is constant in $z$ and
(\ref{eq:pnofa-rho}) reduces to (\ref{eq:pnofa-indep}); this is the
recovery contract that pins the extended model to the inherited one.
Equation~(\ref{eq:pnofa-rho}) requires no multivariate-normal CDF ---
equicorrelation collapses the $\Nloc$-dimensional integral to a single
dimension --- and is the formal content of the rebuild's
``decorrelation channel.''

% ---------------------------------------------------------------------
\section{The two knobs and the question}
\label{sec:knobs}

Equations (\ref{eq:sdt}) expose the observer's two per-location
controls. \emph{Sensitivity} $\dprime_i$ is set by attention through
the asymmetric transfer
$\dprime_c = \dprime_{\max} f(\alpha_c; \beta(r))$,
$\dprime_u = \dprime_{\max} f(\alpha_u; \gamma(r))$, where the
benefit/cost ratio $r$ governs how much sensitivity is gained at the
cued focus versus lost at the periphery. The \emph{criterion} $c_i$
sets the bias independently of sensitivity. Adapting to a valuable cue
can in principle be done by sharpening $\dprime_c$ (moving attention)
or by lowering $c_c$ (shifting the criterion), and the central
question of the paper --- when does value-directed attention matter?
--- is precisely which of these the reward-maximising observer relies
on, and in what parameter regime. With (\ref{eq:pnofa-rho}) promoted to
a model parameter, a third control appears: \emph{decorrelation} of
the cross-location noise, which the inherited model held fixed at
$\rho = 0$.

\begin{thebibliography}{9}
\bibitem{GreenSwets1966}
D.~M. Green and J.~A. Swets.
\emph{Signal Detection Theory and Psychophysics}.
Wiley, New York, 1966.

\bibitem{MacmillanCreelman2005}
N.~A. Macmillan and C.~D. Creelman.
\emph{Detection Theory: A User's Guide}.
2nd ed., Lawrence Erlbaum, 2005.

\bibitem{HermanVDA2026}
Herman Lab.
\emph{When Does Value-Directed Attention Matter? A Normative Model with
Independent Attentional Benefit and Cost}.
Manuscript, 2026-04-09 (target of the rebuild).
\end{thebibliography}

\end{document}
